\section{Explicit bounds and the elementary jet compiler}\label{sec:bounds}

This appendix makes the quantitative analytic steps reproducible. Throughout, $H\ge1$ is an integer and \eqref{eq:Hbounds} holds. All norms of a jet tuple below mean the maximum absolute difference of its finitely many entries.

\subsection{Bounds for the connection and density}

Put $u=b_x-a_y$ and $w=c_x+bP$. Then
\[
|u|,|u_i|\le2H,\qquad
|w|\le3H^3,\qquad |w_i|\le7H^5,
\]
because
\[
P_i=\frac{u_i}{c}-\frac{u c_i}{c^2},\quad
w_i=c_{xi}+b_iP+bP_i,\quad
Q_i=\frac{w_i}{a}-\frac{wa_i}{a^2}.
\]
Successive substitution gives the table
\begin{center}
\begin{tabular}{lc}
\toprule
Quantity & Uniform absolute bound\\
\midrule
$P$ & $2H^2$\\
$P_x,P_y$ & $4H^4$\\
$Q$ & $3H^4$\\
$Q_x,Q_y$ & $10H^6$\\
$\rho=P_y-Q_x$ & $14H^6$\\
$k_g\,\dd s/\dd t$ on any reference side & $13H^4$\\
\bottomrule
\end{tabular}
\end{center}
For the last row, the angle density on $BC$ has numerator
$x\dot y-y\dot x$ with $x=b-a$, $y=c$,
$|x|\le2H$, $|y|\le H$, $|\dot x|\le4H$, $|\dot y|\le2H$.
Its denominator is at least $c^2\ge H^{-2}$, so its absolute value is at most $8H^4$. The $CA$ angle density is at most $2H^4$, and the $AB$ angle density is zero. The connection term is at most $5H^4$ on every side. These calculations prove the uniform bound.

\subsection{Stability with respect to rational jet approximations}

Compare two jet tuples each satisfying the same $H$ bounds, and suppose every entry differs by at most $\eps$. Then
\begin{equation}\label{eq:stabilityconstants}
\begin{aligned}
|\Delta P|&\le4H^3\eps,&
|\Delta P_i|&\le12H^5\eps,\\
|\Delta Q|&\le10H^5\eps,&
|\Delta Q_i|&\le46H^7\eps,\\
|\Delta\rho|&\le58H^7\eps.
\end{aligned}
\end{equation}
Here $\Delta$ means the difference between the two evaluations, not a derivative. The useful elementary estimates are
\[
|c^{-1}-{c'}^{-1}|\le H^2\eps,
\qquad |c^{-2}-{c'}^{-2}|\le2H^3\eps,
\]
and the same for $a$. They follow by factoring the differences, not by a mean-value theorem. Subtract products one factor at a time. This first gives
$|\Delta w|\le7H^4\eps$ and
$|\Delta w_i|\le23H^6\eps$.
The two terms of $Q_i=w_i/a-wa_i/a^2$ contribute at most
$30H^7\eps$ and $16H^7\eps$, respectively. The other rows follow directly from the formulas above.

For the $BC$ angle density, the numerator changes by at most $16H\eps$, the denominator by at most $10H\eps$, and the denominator is at least $H^{-2}$. Factoring a quotient difference gives a bound $96H^7\eps$ for its change. Including the connection term bounds the change of the entire side density by $110H^7\eps$. The other sides have smaller bounds.

For the square coefficients at a fixed rational point $(u,v)$,
\[
|\Delta A|\le14H^5\eps,\quad |\Delta B|\le10H^5\eps,
\]
\[
|\Delta A_v|\le68H^7\eps,\quad
|\Delta B_u|\le102H^7\eps.
\]
Now suppose one tuple is exact and the other comes from evaluation errors at most $\eps\le1/(2H)$. The approximate tuple satisfies the original bounds with $H'=2H$: its positive entries remain at least $1/(2H)$ and all jets are at most $2H$. Replacing $H$ by $2H$ in all the preceding estimates shows that
\begin{equation}\label{eq:universalstability}
|\Delta F|\le32768H^7\eps
\end{equation}
is a common conservative bound for every density or coefficient used in the square-grid proof and the reference-side density. Since
$2^{7\ell(2H)}\ge(2H)^7$, the evaluation precision in \eqref{eq:evalbits} makes \eqref{eq:universalstability} at most
$2^{-(r+m+12)}$, below the required $2^{-(r+8)}/N$.

\subsection{A modulus for the square integrand}

The map $\Phi$ satisfies
$\norm{\Phi(z)-\Phi(z')}_1\le2\norm{z-z'}_1$.
If the square-point distance is at most $\rho_0$ and the resulting jet differences are at most $\eps$, then
\[
\left|\Delta\bigl((1-u)\rho(\Phi)\bigr)\right|
\le58H^7\eps+14H^6\rho_0.
\]
Let $h_H=\ell(H)$ and $s_r=r+8+7h_H$. A sufficient modulus is
\begin{equation}\label{eq:nu}
\nu_\rho(r)=\max\{\mu(s_r)+2,\ s_r+2\}.
\end{equation}
At the prescribed distance, both $\eps$ and $\rho_0$ are at most $2^{-s_r}$, so the displayed bound is below $2^{-r}$. The extra two bits safely cover the $\ell^1$ pullback factor. This is a continuity modulus, not a substitute for the derivative-remainder moduli used in finite differences.

The first derivatives of the square coefficients obey
\[
|A_u|\le28H^6,\quad |A_v|\le17H^6,\quad
|B_u|\le23H^6,\quad |B_v|\le10H^6.
\]
For example, $A_u=P_x-vP_y-vQ_x+v^2Q_y$ after pullback. The corresponding boundary left-sum error is at most
$(2\cdot28+2\cdot10)H^6/N\le80H^6/N$.
The one-dimensional derivative bound implies the needed Lipschitz estimate by the finite FTC, so this argument does not introduce an unproved mean-value principle.

\subsection{Compiling Taylor witnesses by rational operations}

Here is a direct way to generate the moduli $\lambda_A,\lambda_B$ used above. Suppose $f$ has value bound $B_f$, derivative coefficient bound $D_f$, and Taylor remainder bound $\eps\norm h_1$ at a suitably small step. At a step where $\eps\le1$, it also has the increment bound
$|f(z+h)-f(z)|\le C_f\norm h_1$, where $C_f=D_f+1$.

\paragraph{Addition.}
Add the two remainder bounds and request one extra bit for each. Derivative and value bounds add.

\paragraph{Product.}
Writing $\Delta f=f(z+h)-f(z)$ and similarly for $g$,
\[
\Delta(fg)-D(fg)(z)h
=f(z)R_g+g(z)R_f+\Delta f\,\Delta g.
\]
Its absolute value is at most
\begin{equation}\label{eq:productTaylor}
(B_f+B_g)\eps\norm h_1+C_fC_g\norm h_1^2.
\end{equation}
For a target $2^{-r}\norm h_1$, it suffices to use the input moduli at
$r+2+\ell(1+B_f+B_g)$ and to impose
$\norm h_1\le2^{-(r+2+\ell(1+C_fC_g))}$.
Taking the maximum of these finitely many requirements gives a PR output modulus. Bounds may be rounded upward to integers before using $\ell$.

\paragraph{Reciprocal.}
Suppose $|f|\ge1/M$ with a fixed integer $M\ge1$. The exact identity
\[
\frac1{f+\Delta f}-\frac1f
=-\frac{\Delta f}{f^2}
+\frac{(\Delta f)^2}{f^2(f+\Delta f)}
\]
gives the remainder bound
\begin{equation}\label{eq:reciprocalTaylor}
M^2\eps\norm h_1+M^3C_f^2\norm h_1^2.
\end{equation}
Use the input modulus at $r+2+\ell(1+M^2)$ and require
$\norm h_1\le2^{-(r+2+\ell(1+M^3C_f^2))}$.
The lower bound is required at both endpoints, as is already guaranteed on the input domain. This construction works for either sign of $f$.

\paragraph{Fixed compositions.}
For a quantitative map $\psi$ with increment bound $L\norm h_1$ and Taylor remainder $\eps_\psi\norm h_1$, expand
\[
f(\psi(z+h))-f(\psi(z))
=Df(\psi(z))\bigl(\psi(z+h)-\psi(z)\bigr)+R_f.
\]
The remainder after subtracting $Df(\psi(z))D\psi(z)h$ is bounded by
$(L\eps_f+D_f\eps_\psi)\norm h_1$.
Request enough additional bits in both moduli to make this below $2^{-r}\norm h_1$, and ensure $L\norm h_1$ is within $f$'s permitted step size. The polynomial map $\Phi$ has explicit first derivatives and a quadratic remainder, so this is especially simple for the Green proof. Supplied domain margins make all evaluations legal.

Apply the Taylor rules to the fixed expression trees for $P,Q,A,B$. Obtain their first derivatives algebraically, and compile only continuity moduli for those derivative expressions. An input jet of order one has its own derivative contract with the order-two entries as derivatives. Thus the compiler uses only the given $C^2$ jets. It does not ask for third derivatives or for Taylor contracts of $P_i$ and $Q_i$.

\paragraph{Square roots, when starting with metric coefficients.}
For a named $x$ with $1/M\le x\le B$, approximate $x$ by a rational within a prescribed error so the approximation is at least $1/(2M)$. Bisect the rational interval $[0,B+1]$ a predetermined number of times to approximate its square root. Every branch compares rational squares, and the iteration bound is explicit. The identity
\[
|\sqrt x-\sqrt y|=\frac{|x-y|}{\sqrt x+\sqrt y}
\]
supplies a rational Lipschitz bound on this positive range and a PR precision schedule. Differentiation gives
\[
(\sqrt f)_i=\frac{f_i}{2\sqrt f},\qquad
(\sqrt f)_{ij}=\frac{f_{ij}}{2\sqrt f}
-\frac{f_if_j}{4f^{3/2}}.
\]
The same difference identity, followed by the product and reciprocal estimates, proves the Taylor contracts; continuity bounds follow from the displayed formulas. This justifies the effective Cholesky construction without an unbounded search or an assumption that arbitrary square roots are exact rational operations.

The angle evaluator is handled by its explicit integral and the derivative
$(1+x^2)^{-1}$, already proved in \cref{sec:calculus}. Combining these rules supplies every local bound and modulus used in the main precision schedule.
